%! Exponential Functions — Unit 4, Lesson 4.3 — Guided Notes (Answer Key)
\documentclass[10pt]{article}
\usepackage{precalculus-article}
\usepackage{precalculus-key}   % pulls in -boxes; do NOT also load -boxes
\newcommand{\TallMath}[1]{$\displaystyle #1\rule[-1.4em]{0pt}{3.2em}$}

\begin{document}

\pageheader{Unit 4, Lesson 4.3}{Guided Notes --- Answer Key}
\namedateperiod

\vspace{0.06in}

\begin{objectivebox}
By the end of this lesson, I will be able to\ldots
\begin{enumerate}[itemsep=2pt]
  \item Identify the initial value $a$ and base $b$ of $f(x) = ab^x$ and classify as
        \ans{exponential} growth or \ans{exponential} decay. \hfill\textit{(LO 2.3.A.1)}
  \item Describe the \ans{domain}, end \ans{behavior}, \ans{concavity}, and monotonic
        nature of an exponential function. \hfill\textit{(LO 2.3.A.3 / A.5)}
  \item Determine whether a function is exponential by checking whether outputs are
        \ans{proportional} over equal-length input intervals. \hfill\textit{(LO 2.3.A.3 / A.4)}
\end{enumerate}
\end{objectivebox}

\vspace{0.04in}

\begin{vocabbox}
\termblanklong{exponential function}
\ansline{A function $f(x) = ab^x$ with $a \ne 0$, $b > 0$, $b \ne 1$.}
\termblanklong{initial value}
\ansline{The value $a = f(0)$; the $y$-intercept of the function.}
\termblanklong{base}
\ansline{The constant $b$; the multiplicative factor per unit increase in input.}
\termblanklong{exponential growth}
\ansline{When $a > 0$ and $b > 1$; outputs increase without bound.}
\termblanklong{exponential decay}
\ansline{When $a > 0$ and $0 < b < 1$; outputs approach zero.}
\termblanklong{concave up / concave down}
\ansline{Graph curves upward or downward; exponential functions are always one or the other.}
\termblanklong{end behavior}
\ansline{Behavior of $f(x)$ as $x \to \pm\infty$, expressed with limit notation.}
\termblanklong{horizontal asymptote}
\ansline{Line $y = 0$ that the graph of $f(x) = ab^x$ approaches but never reaches.}
\end{vocabbox}

\vspace{0.04in}

\begin{hookbox}
\textbf{The Doubling Penny:} If you receive 1 penny on Day 1, 2 on Day 2, 4 on Day 3, \ldots,
how much total after 30 days?

\medskip
My guess: \textit{(varies)} \quad Actual: $2^{30} - 1 \approx \$\,10.7$ million.

\medskip
The function modeling Day $x$ is $f(x) = 1 \cdot $ \ans{$2$}$^x$.
The initial value is \ans{1} and the base is \ans{2}.
\end{hookbox}

\vspace{0.04in}

\begin{notesbox}{LO 2.3.A.1 --- General Form of an Exponential Function (EK 2.3.A.1)}
\[
  f(x) = a \cdot b^x, \quad a \ne 0,\ b > 0,\ b \ne 1
\]

\begin{itemize}[itemsep=3pt]
  \item $a = f(0) = $ \ans{initial value / $y$-intercept}
  \item $b = $ \ans{common ratio; multiplicative factor per unit input}
  \item \textbf{Exponential growth:} $a > 0$ and $b$ \ans{$>$} 1
  \item \textbf{Exponential decay:} $a > 0$ and \ans{$0$} $< b <$ \ans{$1$}
\end{itemize}

\medskip
\textbf{Example:} $f(x) = 3 \cdot 2^x$. \quad $a = $ \ans{3} \quad $b = $ \ans{2}
\quad Growth or decay? \ans{Growth}

\medskip
\renewcommand{\arraystretch}{1.5}
\begin{tabular}{|c|c|c|c|c|}
\hline
$x$    & 0 & 1 & 2 & 3 \\
\hline
$f(x)$ & \ans{3} & \ans{6} & \ans{12} & \ans{24} \\
\hline
\multicolumn{5}{|l|}{Ratio of successive outputs: \quad $\dfrac{f(1)}{f(0)} = $ \ans{2}
\quad $\dfrac{f(2)}{f(1)} = $ \ans{2} \quad $\dfrac{f(3)}{f(2)} = $ \ans{2}} \\
\hline
\end{tabular}
\end{notesbox}

\vspace{0.04in}

\begin{notesbox}{LO 2.3.A.2/A.3 --- Key Characteristics (EK 2.3.A.2, A.3)}
\renewcommand{\arraystretch}{1.6}
\begin{tabularx}{\linewidth}{@{} >{\bfseries}p{0.28\linewidth} X @{}}
Domain & \ans{All real numbers, $(-\infty, \infty)$} \\
Range (when $a > 0$) & \ans{$(0, \infty)$}; horizontal asymptote: $y = $ \ans{$0$} \\
Monotonicity & Always \ans{increasing} if $b > 1$; always \ans{decreasing} if $0 < b < 1$ \\
Extrema on open intervals & \ans{None} \\
Concavity (when $a > 0$) & Always concave \ans{up}; no \ans{inflection} points \\
\end{tabularx}
\end{notesbox}

\vspace{0.04in}

\begin{notesbox}{LO 2.3.A.5 --- End Behavior (EK 2.3.A.5)}
For $f(x) = ab^x$ with $a > 0$, $b > 1$ (growth):

$\displaystyle\lim_{x \to \infty} f(x) = $ \ans{$\infty$} \qquad
$\displaystyle\lim_{x \to -\infty} f(x) = $ \ans{$0$}

For $f(x) = ab^x$ with $a > 0$, $0 < b < 1$ (decay):

$\displaystyle\lim_{x \to \infty} f(x) = $ \ans{$0$} \qquad
$\displaystyle\lim_{x \to -\infty} f(x) = $ \ans{$\infty$}

\textbf{Example:} State the end behavior of $h(x) = 5 \cdot (0.4)^x$.

\vspace{0.04in}
$\displaystyle\lim_{x\to\infty} h(x) = $ \ans{$0$} \qquad
$\displaystyle\lim_{x\to-\infty} h(x) = $ \ans{$\infty$}
\end{notesbox}

\vspace{0.04in}

\begin{notesbox}{LO 2.3.A.4 --- Detecting Exponential from a Transformed Table (EK 2.3.A.4)}
\textbf{Key idea:} If $g(x) = f(x) + k$ and $f$ is exponential, the outputs of $g$ will
\textit{not} be proportional. To detect $f$: subtract \ans{$k$} from each output of $g$,
then check whether the resulting values have a constant \ans{ratio}.

\medskip
\textbf{Example:} Outputs of $g(x)$: 7, 10, 16, 28 (inputs 0, 1, 2, 3). Suppose $k = 4$.

\vspace{0.04in}
Subtract 4: \ans{3}, \ans{6}, \ans{12}, \ans{24}

Ratios: \ans{2}, \ans{2}, \ans{2}

Are the ratios constant? \ans{Yes} \quad Write $f(x) = $ \ans{$3 \cdot 2^x$}
\quad Write $g(x) = $ \ans{$3 \cdot 2^x + 4$}
\end{notesbox}

\end{document}
